\section{Machine Learning}

\textbf{OBJETIVO GERAL:}

Desenvolver competências para preparar dados, selecionar algoritmos, treinar, validar, otimizar e interpretar modelos de Machine Learning, aplicando metodologias científicas e boas práticas para construção de soluções preditivas em diferentes contextos organizacionais.

\textbf{CARGA HORÁRIA:} 160 horas

\textbf{FUNÇÃO:}

Desenvolver sistemas computacionais baseados em Inteligência Artificial, Análise Preditiva e Ecossistemas IoT, atendendo às normas de qualidade corporativas, robustez metodológica, governança de dados e arquitetura segura.

\textbf{SUBFUNÇÕES:}

\begin{itemize}
\item Preparar conjuntos de dados para treinamento de modelos.
\item Desenvolver modelos supervisionados e não supervisionados.
\item Avaliar e otimizar modelos de Machine Learning.
\item Interpretar resultados e métricas de desempenho.
\item Integrar modelos em aplicações computacionais.
\end{itemize}

\textbf{PADRÕES DE DESEMPENHO:}

\begin{itemize}
\item Realizar análise exploratória dos dados.
\item Preparar dados para treinamento dos modelos.
\item Selecionar algoritmos adequados ao problema.
\item Avaliar desempenho utilizando métricas apropriadas.
\item Ajustar hiperparâmetros.
\item Interpretar resultados obtidos.
\item Aplicar validação cruzada.
\item Desenvolver pipelines de Machine Learning.
\end{itemize}

\textbf{CAPACIDADES TÉCNICAS:}

\textbf{Análise de Dados}

\begin{itemize}
\item Interpretar bases de dados para problemas de Machine Learning.
\item Realizar análise exploratória de dados.
\item Identificar padrões e anomalias.
\item Preparar conjuntos de treinamento e teste.
\end{itemize}

\textbf{Pré-processamento}

\begin{itemize}
\item Tratar dados ausentes.
\item Codificar variáveis categóricas.
\item Normalizar e padronizar dados.
\item Selecionar atributos relevantes.
\end{itemize}

\textbf{Modelagem}

\begin{itemize}
\item Desenvolver modelos supervisionados.
\item Desenvolver modelos não supervisionados.
\item Comparar algoritmos.
\item Selecionar modelos conforme o problema.
\end{itemize}

\textbf{Avaliação}

\begin{itemize}
\item Aplicar métricas de classificação.
\item Aplicar métricas de regressão.
\item Interpretar matrizes de confusão.
\item Realizar validação cruzada.
\end{itemize}

\textbf{Otimização}

\begin{itemize}
\item Ajustar hiperparâmetros.
\item Desenvolver pipelines de Machine Learning.
\item Interpretar importância das variáveis.
\item Aplicar técnicas de explicabilidade.
\end{itemize}

\textbf{CONHECIMENTOS:}

\textbf{Fundamentos de Machine Learning}

\begin{itemize}
\item Aprendizado Supervisionado
\item Aprendizado Não Supervisionado
\item Aprendizado por Reforço (conceitos)
\item Ciclo de Vida de Modelos
\item CRISP-DM
\item CRISP-ML(Q)
\end{itemize}

\textbf{Análise Exploratória de Dados}

\begin{itemize}
\item EDA
\item Visualização de Dados
\item Correlação
\item Distribuições
\item Identificação de Outliers
\end{itemize}

\textbf{Pré-processamento}

\begin{itemize}
\item Limpeza de Dados
\item Normalização
\item Padronização
\item One-Hot Encoding
\item Label Encoding
\item Engenharia de Atributos
\end{itemize}

\textbf{Modelos de Classificação}

\begin{itemize}
\item Regressão Logística
\item K-Nearest Neighbors
\item Naive Bayes
\item Árvores de Decisão
\item Random Forest
\item Support Vector Machines
\end{itemize}

\textbf{Modelos de Regressão}

\begin{itemize}
\item Regressão Linear
\item Regressão Polinomial
\item Regularização Ridge
\item Regularização Lasso
\end{itemize}

\textbf{Modelos Ensemble}

\begin{itemize}
\item Bagging
\item Random Forest
\item Gradient Boosting
\item XGBoost (introdução)
\item LightGBM (introdução)
\end{itemize}

\textbf{Aprendizado Não Supervisionado}

\begin{itemize}
\item K-Means
\item DBSCAN
\item Clusterização Hierárquica
\item PCA
\item Redução de Dimensionalidade
\end{itemize}

\textbf{Avaliação de Modelos}

\begin{itemize}
\item Acurácia
\item Precisão
\item Recall
\item F1-Score
\item ROC-AUC
\item RMSE
\item MAE
\item Matriz de Confusão
\end{itemize}

\textbf{Validação}

\begin{itemize}
\item Hold-out
\item Validação Cruzada
\item Grid Search
\item Random Search
\item Ajuste de Hiperparâmetros
\end{itemize}

\textbf{Interpretabilidade}

\begin{itemize}
\item Feature Importance
\item SHAP (introdução)
\item LIME (introdução)
\item Explicabilidade de Modelos
\end{itemize}

\textbf{CAPACIDADES SOCIOEMOCIONAIS:}

\begin{itemize}
\item Demonstrar pensamento analítico na interpretação dos resultados.
\item Avaliar criticamente o desempenho dos modelos desenvolvidos.
\item Formular hipóteses para melhoria dos modelos.
\item Trabalhar colaborativamente em projetos de Ciência de Dados.
\item Produzir documentação técnica dos experimentos realizados.
\item Comunicar resultados utilizando linguagem técnica adequada.
\item Demonstrar organização durante a condução dos experimentos.
\item Agir com responsabilidade na utilização de modelos preditivos.
\end{itemize}

\textbf{AMBIENTES PEDAGÓGICOS:}

\begin{itemize}
\item Laboratório de Informática.
\item Laboratório de Ciência de Dados.
\item Laboratório de Inteligência Artificial.
\item Ambiente Virtual de Aprendizagem.
\item Ambiente Cloud para treinamento de modelos.
\end{itemize}

\textbf{EQUIPAMENTOS E RECURSOS:}

\textbf{Hardware}

\begin{itemize}
\item Computadores com acesso à Internet.
\item Estações com suporte ao processamento de modelos.
\item Projetor multimídia.
\end{itemize}

\textbf{Software}

\begin{itemize}
\item Python
\item Scikit-Learn
\item Pandas
\item NumPy
\item Matplotlib
\item Seaborn
\item Jupyter Notebook
\item Google Colab
\item VS Code
\item Git
\end{itemize}
